We build a quantitative general equilibrium model that unifies three features
of modern economies absent from standard frameworks: pervasive firm-level
monopsony power in labor markets, heterogeneous households facing incomplete
insurance markets, and multi-bloc trade networks subject to time-varying
geopolitical frictions. Solving via sequence-space Jacobians, we establish
three main results. First, monopsony endogenously flattens the aggregate
Phillips curve---not through menu costs or strategic complementarities, but
through the wage markdown's countercyclical interaction with the earnings
distribution, raising the sacrifice ratio by 40--75\% relative to the
competitive benchmark. Second, supply-side shocks propagate with 30--45\%
greater amplitude and substantially longer persistence in fragmented
input-output networks than in integrated ones, because reshoring concentrates
upstream bottlenecks into fewer, less substitutable domestic nodes with
elevated Domar weights. Third, the standard Taylor rule is welfare-dominated
by a class of augmented rules that condition on measures of labor market
concentration and trade-bloc network entropy, with strictly positive optimal
coefficients on the wage markdown gap and a strictly negative coefficient on
the fragmentation premium. The framework jointly rationalizes the post-2021
inflation surge, the resilience of unemployment to aggressive rate hikes,
and the asymmetric cross-country transmission of energy and supply chain
shocks---three phenomena that standard New Keynesian models cannot
simultaneously explain.

\medskip
\noindent\textbf{JEL Classification:} E12, E24, E52, F41, J42.

\medskip
\noindent\textbf{Keywords:} Monopsony, HANK, input-output networks,
geoeconomic fragmentation, Phillips curve, monetary policy, sequence-space
Jacobians, Domar weights.
